\scanpage{18}
% L01–06
donc loc. \inserted{fppf} ceci provient d'une situation d'un \uline{polygône régulier} \deleted{(projectif)} dans le plan projectif $X/S$, tel que $\deleted{\unsure{\beta}}\ \boxed{\beta\inserted{=\alpha+2}\ \text{inv.}}$ (ceci exprime $a_1\notin H_0$), \deleted{et} \inserted{ou encore tel} que la quadrique \deleted{\unsure{sur}} associée \inserted{$Q$} dans $\unsure{P|X}$ (l'unique quadrique) soit lisse. \deleted{\unsure{car}} Donc

\deleted{\unsure{la réf.} \illegible{1} $\unsure{\frac14|\frac12}\det\left(\begin{smallmatrix}0&&\\&2&\\&&2\end{smallmatrix}\right)$}

% L07–12
\uline{Th} \inserted{$S$ schéma.} a) Soit $C$ la catégorie des couples $(G,\varphi)$ d'une forme $G$ de $\mathrm{GP}(1)_S$ et d'un hom $\varphi:\mathbb D_\infty\longrightarrow G(S)$, \inserted{(d'où $\tau_0,\tau_1,u\in G(S)$)} satisfaisant à
\[
(\ast)\quad\forall s\in S,\qquad \deleted{\unsure{\{}}u_s^2\deleted{\unsure{\}}}\ne1 \qquad
(\text{i.e. }\varphi_s\text{ ne se factorise pas par }\mathbb D_2)
\]
% slanted
\margin{Dire qu'il existe un \uline{unique} div de degré $2$ fixe sous $\mathbb D_\infty$ -- qui correspond à un \unsure{sous-groupe} ``spécial'' de $G$, \unsure{lequel est ce|lequel est un} sous ? \unsure{Cf p.|t.p.}\ldots}

% L13–17
b) $\forall\ \beta\in\Gamma(S,\mathcal O_S^*)$, \deleted{\unsure{soit}} \inserted{soit} $U(\beta)$ le groupe \deleted{\unsure{forme}} lisse \deleted{Soit $C'$ la catégorie des couples formés d'un} de dim $1$ sur $S$, \deleted{défini} formé des $\lambda$ tels que $1+\lambda\beta$ inversible, avec la multiplication
\[
[\lambda]+_{\boldsymbol{.}}[\lambda']=[\lambda+\lambda'+\lambda\lambda'\beta]
\]
Soit $C'$ la catégorie des couples $(\alpha,R)$, où $\alpha\in\Gamma(S,\mathcal O_S)$ \uline{telle que $\beta:\alpha+2$} inv., et $R$ un torseur sous $\deleted{\unsure{G_{2-\alpha}}}\ U(2-\alpha)$.

% L18–23
c) Soit $\alpha\in\Gamma(S,\mathcal O_S)$, \deleted{\unsure{section de} $\mathcal O_S$} \inserted{et considérons} le polygône régulier principal de inv. $\alpha$, d'où, \inserted{si $\alpha+2$ inv. donc $Q_\alpha$ lisse}, une représentation
\[
\varphi_\alpha:\mathbb D_\infty\longrightarrow G_\alpha(S)
\]
\deleted{\unsure{$\ne$} telle \unsure{que}} \inserted{où $G_\alpha=\mathrm{Aut}(X_\alpha,Q_\alpha)$ et} \deleted{$U_{\unsure{\alpha}}$ opère sur $\mathrm{Aut}(G_\alpha,\varphi_\alpha)$.}
\[
U_{2-\alpha}\simeq\uline{\mathrm{Zentr}}(\varphi_\alpha)=\uline{\mathrm{Aut}}(G_\alpha,\varphi_\alpha)
\]
Si $R$ est un torseur sous $G_\alpha$, on peut l'utiliser pour former $(G_\alpha^R,\varphi_\alpha^R)$.

% L24–26
d) On trouve ainsi un foncteur $C'\longrightarrow C^{\unsure{*}}$, et \deleted{\unsure{X}}celui-ci est une équivalence de cat. Le foncteur quasi-inverse est
% archive sheet number: 17
